\section{Conclusion and Outlook}

The central question behind this project work was quite straightforward which was whether a differentiable physics simulator, 
used as the prior mean in Bayesian Optimization, can actually help with a realistic accelerator tuning problem(ARES Experimental Area) 
or does it only work on a toy example(FODO lattice)? Based on everything presented in Chapter 4, it can be confidently answered a “yes”. Under this project scope, the Cheetah based prior idea demonstrated on a simple two dimensional FODO lattice \cite{kaiser2024bridging} was taken and 
scaled it up to the five dimensional beam tuning task at the ARES Experimental Area. The implementation involved building a GPyTorch compatible module 
(\texttt{AresPriorMean}) that wraps the Cheetah simulator and uses six quadrupole misalignment values as trainable parameters. This module is integrated 
with Xopt framework, making it easy to switch between different optimizer setups. \

The results speak for themselves. Both the physics informed BO variants achieved a median MAE of 6$\mu$m, which is significantly higher than standard BO (46$\mu$m) and Nelder-Mead(67$\mu$m).
The matched prior hit the 40$\mu$m target in a median 4 steps, the mismatched prior in 14, whereas standard BO required 82 steps on a rare occasions (20\% success rate) and 
Nelder-Mead never reached the target. What makes the approach especially promising for real world use is its robustness. Even when starting with incorrect misalignments values 
and modified beam distribution (the mismatched prior in BO case), the prior still outperformed the uninformed zero mean prior (standard BO) and caught up to the perfectly matched prior scenario 
of BO very soon within 15 iterations. The trainable misalignment parameters successfully converged very close to the ground truth purely from optimization data, which means that the method does 
not just tune the beam but it also learns something about the machine’s physical state along the way.  \

Ofcourse, the ARES Experimental Ares with its five tuneable parameters is still a relatively small problem compared to large-scale facilities like European XFEL and PETRA III, 
but this is a significant step forward from the two dimensional FODO lattice that was previously tested to evaluate cheetah’s potential. The fact that the method scales gracefully 
to this level sends quite good signs for tackling more complex beamlines in future. 



